% Options for packages loaded elsewhere
\PassOptionsToPackage{unicode}{hyperref}
\PassOptionsToPackage{hyphens}{url}
\PassOptionsToPackage{dvipsnames,svgnames*,x11names*}{xcolor}
%
\documentclass[
  11pt,
  english,
]{article}
\usepackage{lmodern}
\usepackage{amssymb,amsmath}
\usepackage{ifxetex,ifluatex}
\ifnum 0\ifxetex 1\fi\ifluatex 1\fi=0 % if pdftex
  \usepackage[T1]{fontenc}
  \usepackage[utf8]{inputenc}
  \usepackage{textcomp} % provide euro and other symbols
\else % if luatex or xetex
  \usepackage{unicode-math}
  \defaultfontfeatures{Scale=MatchLowercase}
  \defaultfontfeatures[\rmfamily]{Ligatures=TeX,Scale=1}
\fi
% Use upquote if available, for straight quotes in verbatim environments
\IfFileExists{upquote.sty}{\usepackage{upquote}}{}
\IfFileExists{microtype.sty}{% use microtype if available
  \usepackage[]{microtype}
  \UseMicrotypeSet[protrusion]{basicmath} % disable protrusion for tt fonts
}{}
\makeatletter
\@ifundefined{KOMAClassName}{% if non-KOMA class
  \IfFileExists{parskip.sty}{%
    \usepackage{parskip}
  }{% else
    \setlength{\parindent}{0pt}
    \setlength{\parskip}{6pt plus 2pt minus 1pt}}
}{% if KOMA class
  \KOMAoptions{parskip=half}}
\makeatother
\usepackage{xcolor}
\IfFileExists{xurl.sty}{\usepackage{xurl}}{} % add URL line breaks if available
\IfFileExists{bookmark.sty}{\usepackage{bookmark}}{\usepackage{hyperref}}
\hypersetup{
  pdftitle={Humanvoice},
  pdfauthor={Product proposal},
  pdflang={en},
  colorlinks=true,
  linkcolor=blue,
  filecolor=Maroon,
  citecolor=Blue,
  urlcolor=blue,
  pdfcreator={LaTeX via pandoc}}
\urlstyle{same} % disable monospaced font for URLs
\usepackage[margin=1.05in]{geometry}
\usepackage{longtable,booktabs}
% Correct order of tables after \paragraph or \subparagraph
\usepackage{etoolbox}
\makeatletter
\patchcmd\longtable{\par}{\if@noskipsec\mbox{}\fi\par}{}{}
\makeatother
% Allow footnotes in longtable head/foot
\IfFileExists{footnotehyper.sty}{\usepackage{footnotehyper}}{\usepackage{footnote}}
\makesavenoteenv{longtable}
\setlength{\emergencystretch}{3em} % prevent overfull lines
\providecommand{\tightlist}{%
  \setlength{\itemsep}{0pt}\setlength{\parskip}{0pt}}
\setcounter{secnumdepth}{-\maxdimen} % remove section numbering
\ifxetex
  % Load polyglossia as late as possible: uses bidi with RTL langages (e.g. Hebrew, Arabic)
  \usepackage{polyglossia}
  \setmainlanguage[]{english}
\else
  \usepackage[shorthands=off,main=english]{babel}
\fi
\usepackage[]{natbib}
\bibliographystyle{plainnat}

\title{Humanvoice}
\usepackage{etoolbox}
\makeatletter
\providecommand{\subtitle}[1]{% add subtitle to \maketitle
  \apptocmd{\@title}{\par {\large #1 \par}}{}{}
}
\makeatother
\subtitle{A reader-first revision system for technical writing}
\author{Product proposal}
\date{23 August 2026}

\begin{document}
\maketitle
\begin{abstract}
Technical authors now have abundant help generating and polishing prose,
but they still lack a reliable way to move a complex document from
locally fluent to understood and accepted. Humanvoice is a revision
system for that gap. It locates diagnosable problems without rewriting
automatically, protects equations, numbers, citations, and claims during
revision, and ends with a decision from a named human reader. This
proposal requests a bounded twelve-week MVP and one-document feasibility
run. The build will establish whether the workflow is usable and
measurable; it will not be presented as evidence that the product
improves expert writing. A later comparative study must earn that claim.
\end{abstract}

\hypertarget{decision-in-brief}{%
\section{Decision in brief}\label{decision-in-brief}}

Build humanvoice as a reader-first revision system for technical
documents. The product should help an author find where a draft loses
its reader, revise without silently changing its substance, and learn
whether the next version actually reduced the work required for
acceptance.

The immediate decision is deliberately bounded. Authorize a twelve-week
MVP, an evaluation corpus, and one live-document feasibility run. At the
end of that work, decide whether to fund a comparative study against the
tools and review process already in use. Do not authorize general
deployment, and do not claim that humanvoice improves writing before
that comparison is complete.

The MVP can be built by one product engineer and one research and
evaluation lead, with a part-time domain reviewer. The primary outcome
for the eventual study is the number of human-feedback rounds required
to reach reader acceptance, provided that reader time and
meaning-preservation failures do not increase. If the product cannot
improve that outcome, it has not earned a place in the writing process.

\hypertarget{the-problem-worth-solving}{%
\section{The problem worth solving}\label{the-problem-worth-solving}}

A technical document can be correct, compile without warnings, and
survive several agent reviews while failing in its first few minutes
with a human reader. One funding proposal in this workspace did exactly
that. Its reader could not state the purpose, did not know why internal
file names appeared in the argument, and could not tell what action was
being requested. The document had not failed because of a typo. It had
transferred the work of constructing the argument from the author to a
scarce expert reader.\footnote{The full case record, including the
  reader's response and the later repair history, is preserved in
  \href{part1_failures.tex}{the survey's internal evidence appendix}. It
  is used here as a motivating case, not as external efficacy evidence.}

That transfer has an economic cost. A rejected draft begins another
cycle of diagnosis, revision, and rereading. Senior readers spend time
rediscovering the same kinds of problems: private project language in
public exposition, repetitive paragraph structure, qualifications that
never explain a mechanism, unsupported specificity, and missing links
between an equation and the decision it is supposed to inform. The cost
is not simply author time. Every cycle consumes the attention of the
person whose judgment the document was written to support.

Current tools make only parts of this cycle cheaper. Grammar checkers
and prose linters find local surface defects. General language models
can draft, rewrite, and comment. Citation services can confirm that a
DOI exists. Version control can show that text changed. None of these
tells us that the intended reader can now reconstruct the argument, and
none protects every substantive object while a large passage is
recomposed.

Humanvoice is aimed at this missing middle: the journey from a plausible
draft to a document a named reader understands and accepts.

\hypertarget{why-the-obvious-alternative-stops-short}{%
\section{Why the obvious alternative stops
short}\label{why-the-obvious-alternative-stops-short}}

The strongest rival is not another specialist writing product. It is the
combination most authors can assemble today: an ordinary prose linter, a
general LLM editor, version control, and manual review. Humanvoice
should reuse those components where they work and compete on the
connection between them.

\begin{longtable}[]{@{}lll@{}}
\toprule
\begin{minipage}[b]{0.30\columnwidth}\raggedright
Revision job\strut
\end{minipage} & \begin{minipage}[b]{0.30\columnwidth}\raggedright
Ordinary toolchain\strut
\end{minipage} & \begin{minipage}[b]{0.30\columnwidth}\raggedright
Humanvoice\strut
\end{minipage}\tabularnewline
\midrule
\endhead
\begin{minipage}[t]{0.30\columnwidth}\raggedright
Find local style and grammar problems\strut
\end{minipage} & \begin{minipage}[t]{0.30\columnwidth}\raggedright
Strong\strut
\end{minipage} & \begin{minipage}[t]{0.30\columnwidth}\raggedright
Reuses replaceable linters and adds project-register diagnostics\strut
\end{minipage}\tabularnewline
\begin{minipage}[t]{0.30\columnwidth}\raggedright
Explain which issue matters most\strut
\end{minipage} & \begin{minipage}[t]{0.30\columnwidth}\raggedright
Inconsistent, often presented as fluent commentary\strut
\end{minipage} & \begin{minipage}[t]{0.30\columnwidth}\raggedright
Ranks located findings and leaves the editorial decision visible\strut
\end{minipage}\tabularnewline
\begin{minipage}[t]{0.30\columnwidth}\raggedright
Protect equations, numbers, citations, and claims\strut
\end{minipage} & \begin{minipage}[t]{0.30\columnwidth}\raggedright
Depends on manual comparison\strut
\end{minipage} & \begin{minipage}[t]{0.30\columnwidth}\raggedright
Treats them as protected content and requires an explanation for
substantive change\strut
\end{minipage}\tabularnewline
\begin{minipage}[t]{0.30\columnwidth}\raggedright
Determine whether the revision helped\strut
\end{minipage} & \begin{minipage}[t]{0.30\columnwidth}\raggedright
Surface scores or an uncalibrated model judgment\strut
\end{minipage} & \begin{minipage}[t]{0.30\columnwidth}\raggedright
Compares reader reconstruction, review burden, and acceptance\strut
\end{minipage}\tabularnewline
\begin{minipage}[t]{0.30\columnwidth}\raggedright
Learn across documents\strut
\end{minipage} & \begin{minipage}[t]{0.30\columnwidth}\raggedright
Prompt history and scattered review notes\strut
\end{minipage} & \begin{minipage}[t]{0.30\columnwidth}\raggedright
Retains bounded, privacy-aware revision and reader outcomes\strut
\end{minipage}\tabularnewline
\bottomrule
\end{longtable}

The distinction matters because the evidence for AI assistance is
conditional. In a preregistered experiment, generative assistance
reduced completion time and improved evaluated output on
occupation-specific writing tasks \citep{noy2023experimental}. Field
evidence also reports time or productivity gains in workplace settings
\citep{dillon2025shifting, brynjolfsson2025work}. These studies
establish that assistance can change work. They do not estimate the
value of this workflow for expert economics writing, and the average
effects hide differences by task and worker experience. Assistance can
also reduce correctness when a task lies outside the model's capability
frontier \citep{dellacqua2026jagged}.

The appropriate product claim is therefore modest and testable:
humanvoice may reduce wasted revision and rereading by connecting
diagnosis, protected editing, and reader evidence. The feasibility build
is designed to learn whether that mechanism can work. It is not a
demonstration that it already does.

\hypertarget{the-product-in-use}{%
\section{The product in use}\label{the-product-in-use}}

Humanvoice sits beside the author's existing editor and repository. It
does not replace either. A document moves through four observable steps.

First, the author opens a LaTeX or Markdown draft and chooses the
intended reader and decision. Humanvoice parses the document into
ordinary prose and protected content. Equations, numbers, labels,
citations, quotations, and declared claims retain their source
locations.

Second, the system presents a short, ranked set of findings. Some are
fully deterministic: a private path appears in the manuscript, five
consecutive sections open the same way, a citation has mismatched
metadata, or a number changed between versions. Others are editorial
questions attached to a span: does this qualification prevent a real
misunderstanding, or does it merely protect the author? Is the paragraph
missing the mechanism that makes its conclusion follow? Findings remain
findings. Humanvoice does not silently edit the source.

Third, the author or an explicitly chosen editor revises the draft. A
protected comparison explains what changed and calls attention to any
altered equation, number, citation, quotation, or claim. The author may
make a substantive change, but it cannot disappear inside a smoother
sentence.

Finally, a low-context reviewer receives the rendered document without
the repository history that made its private language seem familiar. The
reviewer answers questions tied to the document's purpose: What problem
is being solved? What is proposed? Why is the obvious alternative
insufficient? What evidence matters? What decision is requested? A named
human remains the final authority. A missing answer sends the document
back for substantive revision, not another round of synonym replacement.

\hypertarget{revision-flow}{%
\subsection{Revision flow}\label{revision-flow}}

\begin{verbatim}
source -> format-aware parse -> located findings -> author revises
                 |                                  |
                 +-> protected content -------------+
                                                    |
                                           protected comparison
                                                    |
                                   low-context read -> human decision
\end{verbatim}

This design keeps two records for good reason. The reader sees only the
document. The project retains the minimum provenance needed to reproduce
a diagnostic, understand a revision, and evaluate the workflow. Internal
IDs, search administration, and review machinery never become manuscript
prose.

\hypertarget{what-the-evidence-changes}{%
\section{What the evidence changes}\label{what-the-evidence-changes}}

The literature does not point toward an autonomous rewriting system. It
points toward a bounded editor with direct evaluation of the resulting
revision.

\hypertarget{feedback-must-be-judged-by-the-next-draft}{%
\subsection{Feedback must be judged by the next
draft}\label{feedback-must-be-judged-by-the-next-draft}}

Writing feedback can sound specific while missing the most important
problem. On a test set of 1,300 deliberately corrupted stories, models
often produced accurate local comments but failed to identify the
central defect or choose appropriately between criticism and praise
\citep{rashkin2025story}. Work on an economic essay task similarly
evaluates feedback through the revisions it induces rather than through
the fluency of the comment itself \citep{nair2024closing}. Humanvoice
therefore links each finding to a source span, records whether the
writer acted on it, and evaluates what changed in the next version.

Two recent systematic reviews reinforce the boundary. Across their
respective evidence bases, benefits are more consistent for grammar,
process, coherence, and surface accuracy than for argumentation,
discourse organization, or creativity
\citep{meng2026systematic, abudalfa2026awe}. Humanvoice can automate
some measurements at the surface. It must measure higher-order
understanding with the intended readers.

\hypertarget{model-judgments-need-calibration}{%
\subsection{Model judgments need
calibration}\label{model-judgments-need-calibration}}

Model judges are useful screening instruments, but their preferences
change with answer order, length, and model family
\citep{zheng2023judging, panickssery2024llm, dubois2024length}. They
also tend to be more generous than expert writers when judging creative
prose \citep{chakrabarty2024art}. Any model-assisted review in
humanvoice will therefore address one construct at a time, randomize
pair order, control for length, use a different model family from the
editor, and report calibration against human labels. It may direct
attention. It may not accept a document.

\hypertarget{assistance-can-change-trust-and-voice}{%
\subsection{Assistance can change trust and
voice}\label{assistance-can-change-trust-and-voice}}

Disclosure of AI assistance can affect reader ratings and increase
disagreement among readers \citep{li2024disclosure}. Assistance can also
increase individual performance while reducing the diversity of a
collection \citep{doshi2024generative, padmakumar2024writing}. These are
product outcomes, not matters of cosmetic style. Humanvoice will retain
the disclosure condition used in evaluation and monitor convergence
across a sufficiently large corpus. It will never infer authorship from
an individual document or optimize text to evade a detector.

\hypertarget{citation-identity-is-not-claim-support}{%
\subsection{Citation identity is not claim
support}\label{citation-identity-is-not-claim-support}}

Generated scholarly references can be nonexistent or mismatched
\citep{mugaanyi2024citation}. A metadata service can establish that a
paper and DOI match. Only inspection of the paper can establish that it
supports the sentence in which it is cited. Humanvoice separates those
jobs: automatic identity resolution followed by a human source-to-claim
check for important claims.

These findings narrow the design and the promise. They justify local
measurement, protected revision, explicit uncertainty, and direct reader
outcomes. They do not justify a single quality score, an automatic
rewrite, an authorship detector, or an efficacy claim before the pilot.

\hypertarget{what-we-will-build}{%
\section{What we will build}\label{what-we-will-build}}

The MVP has four product capabilities and one evaluation asset.

The document core reads LaTeX and Markdown while preserving source
locations and protected regions. Pandoc and pylatexenc are the first
parser candidates because they expose structured documents through
mature, documented interfaces \citep{pandoc2026filters, pylatexenc2026}.
They will be tested on unknown macros, nested environments, comments,
mathematics, citations, malformed input, and round-trip preservation. A
parser is selected only after those fixtures pass.

The diagnostic layer locates private project register, structural
repetition, defensive prose candidates, citation identity problems, and
build or rendering defects. Vale and existing local rules are plausible
hosts for deterministic prose checks \citep{vale2026docs}. Every result
carries its source span and a clear question or measurement. No
diagnostic command mutates its input.

The protected comparison reports changed equations, numbers, labels,
citations, quotations, and declared claims. latexdiff may provide a
useful rendered comparison, but the invariant check remains structural
\citep{latexdiff2026}. An unexplained protected change stops that
revision from being presented to a reader.

The review workspace shows a clean rendered draft, asks a short
document- specific rubric, and records the minimum decision evidence. It
does not expose the author's plans or audit records. Model-generated
suggestions are optional, identified, and reversible.

The evaluation corpus begins with existing rejected and repaired
documents for which reader verdicts or revision histories are already
available. It will add adjudicated labels for defect location, issue
importance, protected changes, and acceptance. Tuning and held-out
documents will be separated before package or model results are
reported.

The detailed command contracts and sixteen implementation requirements
remain available in the
\href{humanvoice_product_requirements.csv}{requirements trace}. They are
engineering inputs, not the public explanation of the product.

\hypertarget{how-we-will-know-whether-it-works}{%
\section{How we will know whether it
works}\label{how-we-will-know-whether-it-works}}

The evaluation has two stages because feasibility and efficacy answer
different questions.

The first live document tests whether the pieces work together. Can the
parser protect the document? Do findings point to the right source
spans? Can an author understand and dismiss a false positive? Does the
comparison expose meaning changes? Can a reader complete the rubric
without seeing private project machinery? Does the workflow respect
privacy and disclosure choices? One document can reveal integration
failures and unacceptable burden. It cannot estimate an effect.

If feasibility passes, the project will preregister a paired or stepped
comparison against the current workflow: existing linters, ordinary LLM
editing where the author already uses it, version control, and manual
review. Documents will be selected before outcomes are known. Reader
blinding will be used where practical. The sample size will be chosen by
simulating power or precision from the observed distribution of baseline
revision rounds and reader clustering, not by choosing a convenient
number of documents.

The primary estimand is

\[
\Delta r = E[r_{\text{humanvoice}} - r_{\text{current}}],
\]

where \(r\) is the number of human-feedback rounds to acceptance.
Deployment requires evidence that \(\Delta r < 0\) without an increase
in meaning- preservation failures or total reader time. Reader
reconstruction accuracy, substantive revision rate, false-positive
burden, editing time, suggestion acceptance, and cross-document
convergence are explanatory outcomes. None may replace the primary
endpoint.

Results will be reported separately by writer experience, document
genre, and reader role. This is not optional subgroup decoration. Field
evidence shows that average assistance effects can conceal smaller gains
or harm among the most experienced workers \citep{brynjolfsson2025work}.

The project stops or changes direction when any of the following occurs:

\begin{itemize}
\tightlist
\item
  unexplained changes to protected content exceed the prespecified
  tolerance;
\item
  diagnostic false positives consume more review time than they save;
\item
  privacy or disclosure requirements cannot be met;
\item
  a feasibility user cannot understand or control the workflow; or
\item
  reader acceptance does not improve, or becomes slower, after two
  comparative pilot cycles.
\end{itemize}

\hypertarget{delivery-and-investment}{%
\section{Delivery and investment}\label{delivery-and-investment}}

The twelve-week program ends with a working feasibility system, not a
claim of market readiness.

\begin{longtable}[]{@{}rll@{}}
\toprule
\begin{minipage}[b]{0.37\columnwidth}\raggedleft
Weeks\strut
\end{minipage} & \begin{minipage}[b]{0.27\columnwidth}\raggedright
Work\strut
\end{minipage} & \begin{minipage}[b]{0.27\columnwidth}\raggedright
Reader-visible result\strut
\end{minipage}\tabularnewline
\midrule
\endhead
\begin{minipage}[t]{0.37\columnwidth}\raggedleft
1--2\strut
\end{minipage} & \begin{minipage}[t]{0.27\columnwidth}\raggedright
Freeze the baseline corpus, reader rubric, protected-content model, and
parser fixtures\strut
\end{minipage} & \begin{minipage}[t]{0.27\columnwidth}\raggedright
One representative document parses and renders with protected spans
intact\strut
\end{minipage}\tabularnewline
\begin{minipage}[t]{0.37\columnwidth}\raggedleft
2--5\strut
\end{minipage} & \begin{minipage}[t]{0.27\columnwidth}\raggedright
Implement located diagnostics, citation identity checks, and protected
comparison\strut
\end{minipage} & \begin{minipage}[t]{0.27\columnwidth}\raggedright
An author can inspect findings and revise without hidden source
mutation\strut
\end{minipage}\tabularnewline
\begin{minipage}[t]{0.37\columnwidth}\raggedleft
3--7\strut
\end{minipage} & \begin{minipage}[t]{0.27\columnwidth}\raggedright
Label tuning and held-out documents; benchmark parser and diagnostic
candidates\strut
\end{minipage} & \begin{minipage}[t]{0.27\columnwidth}\raggedright
Package decisions are supported by observed behavior rather than
popularity\strut
\end{minipage}\tabularnewline
\begin{minipage}[t]{0.37\columnwidth}\raggedleft
6--9\strut
\end{minipage} & \begin{minipage}[t]{0.27\columnwidth}\raggedright
Integrate the author and reader workflow with local-first data
handling\strut
\end{minipage} & \begin{minipage}[t]{0.27\columnwidth}\raggedright
One document can move from source through revision to a recorded reader
response\strut
\end{minipage}\tabularnewline
\begin{minipage}[t]{0.37\columnwidth}\raggedleft
9--12\strut
\end{minipage} & \begin{minipage}[t]{0.27\columnwidth}\raggedright
Run the live feasibility document, measure burden and failures, and
price the next study\strut
\end{minipage} & \begin{minipage}[t]{0.27\columnwidth}\raggedright
Proceed, revise, or stop recommendation with a preregisterable efficacy
design\strut
\end{minipage}\tabularnewline
\bottomrule
\end{longtable}

The planning assumption is one product engineer, one research and
evaluation lead, and a part-time economics or scholarly-writing
reviewer. Twelve weeks covers the MVP and feasibility case. The efficacy
study's duration depends on the required precision and the arrival rate
of suitable live documents.

The investment case should be calculated from project data. Let \(r_0\)
and \(r_1\) be revision rounds under the current and proposed workflows,
\(t_0\) and \(t_1\) the expert-reader hours per round, \(c_r\) the
loaded cost of that time, \(N\) the annual number of relevant documents,
and \(C\) the annualized build and maintenance cost. The measurable
reader-time component of annual value is

\[
V = N(r_0t_0-r_1t_1)c_r-C.
\]

Quality, author time, and avoided meaning failures may add value, but
they should enter the decision only when the pilot measures them. No
published study supplies the project-specific inputs to this equation.

\hypertarget{risks-and-limits}{%
\section{Risks and limits}\label{risks-and-limits}}

The most serious product risk is meaning drift. Smoother prose is a loss
if it changes a model, number, source, or conclusion. Protected
comparison and human source checks are therefore part of the core, not
later compliance features.

A second risk is alert fatigue. A diagnostic that is technically correct
but rarely useful can make the workflow slower and teach authors to
ignore it. The held-out benchmark will report precision by category and
the time required to review findings. Low-value rules will be removed
rather than defended by their coverage.

A third risk is convergence on a house style. Humanvoice is not meant to
make every economist sound alike. It reports patterns and asks
questions; it does not automatically normalize wording. Aggregate
monitoring will look for loss of lexical and structural diversity, while
individual documents will never receive authorship labels.

Privacy and disclosure create a fourth boundary. Protected manuscripts
remain local by default. Any remote model use must identify the
provider, purpose, retention terms, and approved content. The reader
study must use the disclosure condition intended for actual use because
disclosure itself can affect reader response.

Finally, the supporting evidence is not yet a completed systematic
review. The current audit has broad public discovery and verified
bibliographic identities, but independent screening, specialist database
coverage, full-text verification, critical appraisal, held-out package
benchmarks, and external review remain incomplete. The detailed status
is reported in the \href{humanvoice_product_proposal_audit.md}{evidence
audit}. This limitation blocks strong literature and efficacy claims. It
does not prevent a reversible build whose purpose is to measure
feasibility.

\hypertarget{decision-requested}{%
\section{Decision requested}\label{decision-requested}}

Authorize the twelve-week MVP and one live-document feasibility run
under the resource assumptions and stop rules above. The authorized
deliverables are:

\begin{enumerate}
\def\labelenumi{\arabic{enumi}.}
\tightlist
\item
  a local-first revision system for LaTeX and Markdown;
\item
  located, non-mutating diagnostics and protected comparison;
\item
  a small adjudicated corpus with frozen tuning and held-out partitions;
\item
  one end-to-end feasibility result, including burden and failure
  evidence;
\item
  a costed, power-justified design for the comparative study; and
\item
  a recommendation to proceed, revise, or stop.
\end{enumerate}

This decision does not authorize autonomous rewriting, authorship
detection, general deployment, or a claim that humanvoice improves
expert writing. It authorizes the smallest build that can determine
whether those further investments are warranted.

\hypertarget{technical-appendix}{%
\section{Technical appendix}\label{technical-appendix}}

\hypertarget{product-boundary}{%
\subsection{Product boundary}\label{product-boundary}}

The MVP accepts local LaTeX and Markdown sources and emits findings,
protected comparisons, rendered review material, and bounded evaluation
events. It does not generate claims or citations, classify authorship,
evade detectors, upload manuscripts without approval, or accept a
document on behalf of its reader.

The stable CLI concepts from the founding survey remain useful
implementation interfaces:

\begin{longtable}[]{@{}ll@{}}
\toprule
\begin{minipage}[b]{0.47\columnwidth}\raggedright
Command\strut
\end{minipage} & \begin{minipage}[b]{0.47\columnwidth}\raggedright
Contract\strut
\end{minipage}\tabularnewline
\midrule
\endhead
\begin{minipage}[t]{0.47\columnwidth}\raggedright
\texttt{hv\ leak}\strut
\end{minipage} & \begin{minipage}[t]{0.47\columnwidth}\raggedright
Locate private project register and disclosure problems without
editing\strut
\end{minipage}\tabularnewline
\begin{minipage}[t]{0.47\columnwidth}\raggedright
\texttt{hv\ rhythm}\strut
\end{minipage} & \begin{minipage}[t]{0.47\columnwidth}\raggedright
Report sentence, paragraph, and section-shape measurements against genre
configuration\strut
\end{minipage}\tabularnewline
\begin{minipage}[t]{0.47\columnwidth}\raggedright
\texttt{hv\ defensive}\strut
\end{minipage} & \begin{minipage}[t]{0.47\columnwidth}\raggedright
Return contextualized candidate passages and editorial questions\strut
\end{minipage}\tabularnewline
\begin{minipage}[t]{0.47\columnwidth}\raggedright
\texttt{hv\ diff}\strut
\end{minipage} & \begin{minipage}[t]{0.47\columnwidth}\raggedright
Explain changes to protected content between document versions\strut
\end{minipage}\tabularnewline
\begin{minipage}[t]{0.47\columnwidth}\raggedright
\texttt{hv\ cite}\strut
\end{minipage} & \begin{minipage}[t]{0.47\columnwidth}\raggedright
Resolve bibliographic identity and prepare a separate source-to-claim
review queue\strut
\end{minipage}\tabularnewline
\begin{minipage}[t]{0.47\columnwidth}\raggedright
\texttt{hv\ gate}\strut
\end{minipage} & \begin{minipage}[t]{0.47\columnwidth}\raggedright
Build and render the document, report reference and layout failures, and
expose pages for inspection\strut
\end{minipage}\tabularnewline
\begin{minipage}[t]{0.47\columnwidth}\raggedright
\texttt{hv\ drift}\strut
\end{minipage} & \begin{minipage}[t]{0.47\columnwidth}\raggedright
Measure population-level convergence on sufficiently large corpora;
refuse individual attribution\strut
\end{minipage}\tabularnewline
\bottomrule
\end{longtable}

The names are less important than the shared contracts: immutable input,
stable source locations, structured output, no silent source mutation,
and no diagnostic promotion of a document.

\hypertarget{core-records}{%
\subsection{Core records}\label{core-records}}

Implementation needs six small records: a document version with
protected spans; a located diagnostic finding; an explained change; a
citation identity and claim anchor; a reader response tied to one
rendered version; and a privacy-bounded interaction event. Their
complete fields and evidence links are maintained in the
machine-readable requirements rather than repeated in the main proposal.

\hypertarget{supporting-scholarship-and-audit}{%
\subsection{Supporting scholarship and
audit}\label{supporting-scholarship-and-audit}}

The full literature and software treatment is in
\href{humanvoice_survey.pdf}{the supporting survey}. Search histories,
screening queues, evidence extraction, package inventory, manifests, and
unresolved review work remain under
\href{audit/latest/}{\texttt{docs/survey/audit/latest/}}. Those records
protect the product's claims and implementation choices. They are not
the product proposal's narrative.

  \bibliography{humanvoice_survey}

\end{document}
